\noindent
This review surveys the measurement technology behind commercial golf launch
monitors --- Doppler-radar systems (TrackMan, FlightScope, Garmin~R10, Full
Swing KIT), photometric camera systems (Foresight GCQuad/GC3, Uneekor,
SkyTrak, ProTee~VX, Garmin~R50), and the hybrid radar+camera architectures
that now dominate the high end --- with emphasis on \emph{how each system
measures or infers club-delivery parameters} (club speed, club path, attack
angle, face angle, dynamic loft, impact location) \emph{and ball-launch
parameters} (ball speed, launch angles, spin rate, spin axis). Primary
sources include the governing patent families (TrackMan/Tuxen radar spin and
spin-axis patents, the FlightScope/EDH dielectric-lens spin patent, the
Wintriss/Foresight photometric patents, and the Acushnet stereo
launch-monitor lineage), regulatory (FCC) filings, manufacturer technical
documentation, and the peer-reviewed validation literature. The review
develops the underlying physics --- the D-plane impact model, oblique-impact
spin generation, gear effect, and golf-ball aerodynamics --- as the common
mathematical core every launch monitor implements, and closes with a
freedom-to-operate map of the patent landscape and concrete architectural
recommendations for OpenFlight, an open-source Doppler-radar launch monitor
built on the OmniPreSense OPS243-A and low-cost angle-radar modules. The
central finding: every parameter a launch monitor reports sits on a
\emph{directness hierarchy} from measured to modeled, the industry has
converged on optical augmentation of radar (and radar augmentation of
cameras) precisely because neither modality can measure the full parameter
set alone, and the recently expired Wintriss photometric patents
(2025) open a practical low-cost optical path to true club-face and spin
measurement that complements OpenFlight's radar core.
